% !TeX root = ../../../research-log.tex

\subsection*{dual Hopf algebra}

Let $k$ be a field, and let $(H, \nabla, \eta, \delta, \epsilon, S)$ be a Hopf algebra over $k$.

Following \cite{DNR}, we define the finite dual of $H$.
\[
H^\circ := \left\{ f \in H^* \middle| \ \exists \text{ a two-sided ideal } I \text{ in } \ker(f) \text{ such that } \dim_k(A/I) < \infty \right\}
\]

\begin{proposition}
    Let $f \in H^*$.
    The followings are equivalent

    \begin{enumerate}
        \item There exists a two-sided ideal $I \subset \ker(f)$ of $A$ such that $\dim(A/I) < \infty$.
        \item $\nabla^*(f)\in H^* \otimes H^*$.
        \item There exists finitely many $g_i, h_i \in H^*$ such that, for all $a, b \in H$, $f(ab) = \sum g_i(a) h_i(b)$.
        \item $\operatorname{span} \left\{ L_a(f) \middle| ^\forall a \in H \right\}$ is finite dimensional vector space.
        \item $f$ is a matrix coeffciient of a finitely dimensional $H$-module.
    \end{enumerate}

\end{proposition}
\begin{proof}
    Let write a bilinear map $B_f := \nabla^*(f): H \otimes H \rightarrow k$, and a linear map $T_f :H \rightarrow H^*$ by $T_f(a) := L_a(f)$ for all $a \in H$. 
    Denote $\ker(T_f)$ by $J$, then it is easy to verify that $J$ is a right ideal.
    If we define $I := \left\{ a \in H \middle| Aa \subseteq J \right\} \subset J$ then $I$ is a two-sided ideal.
    \begin{itemize}
        \item $ 2 \Leftrightarrow 3$: As $B_f(a \otimes b) = f(ab)$, it is just for writing out the definition.
        \item $ 3 \Leftrightarrow 4$: If we assume $(3)$ then $T_f(a) = \sum g_i(a) h_i$, and hence $\operatorname{Im}T_f < \infty$.
        Conversely, if one take a basis $\{ h_i \}$ of $\operatorname{Im}T_f$ and the dual basis $\{\lambda^i \}$ of its dual, then 
        \[
        T_f = \sum \lambda^i (T_f) h_i
        \]
        and, if we define $g_i := \lambda^i(T_f)$, then $T_f = \sum g_i \otimes h_i \in A^* \otimes A^*$.
        \item $ 1 \Leftrightarrow 4$: Since $B_f$ factors through $H/I \otimes H/I$, $\operatorname{Im}T_f$ is of finite dimensional which means that $B_f$ is of finite rank.
        Conversely, note that $\operatorname{Im}T_f \cong H/J < \infty$, and $I$ is the kernel of
        \[
        H^{\text{op}} \rightarrow \End_k(H/J), \quad b \mapsto (a+J \mapsto ab+J).
        \]
        Hence, there exsits an embedding $H/I \hookrightarrow \End_k(H/J)$ and it implies that $H/I$ is a finite dimensional vector space.
        \item $ 5 \Rightarrow 1$: Let $V$ be a finite dimensional $H$-module, and for $v \in V, \lambda \in V^*$, let define $c_{\lambda, v}(a) := \lambda(a \cdot v)$.
        If we define $\rho: H \rightarrow \End_k(V)$ and then $\ker(\rho)$ is a two-sided ideal with finite codimension inside in $\ker(c_{\lambda, v})$.
        \item $ 5 \Rightarrow 2$:
        \begin{align*}
            c_{\lambda, v}(Ab) & = \lambda(ab \cdot v) \\
            & = \lambda_j e^j (\rho(a) e^i(\rho(b)v)e_i) \\
            & = \lambda_j e^j(\rho(a)e_i) \cdot e^i(\rho(b)v).
        \end{align*}
        By defining $g_i := \lambda_j e^j(\rho(a)e_i)$ and $h_i := e^i(\rho(b)v)$ then $c_{\lambda,v} = \sum g_i h_i$.
        \item $ 1 \Rightarrow 5$ If we let $V = H / I$ and let $\lambda(a+I):=f(a)$ and $v = 1 + I$, then we have
        \[
            f(a) = f(a + I) = \lambda(a \cdot v)
        \]
        which implies that $f$ is a matrix coefficient.
    \end{itemize}
\end{proof}

Also, $(-)^\circ$ s a functorial; $(-)^\circ$ is an endofunctor on $\mathrm{Hopf}_k$.
Since, for any Hopf algebra morphism $\phi$, it takes a matrix coefficient of a module to a matrix coefficient of the same module but over differnet Hopf algebra.


\subsection*{moddules and comodules of $A$, $A^*$ and $A^\circ$.}

Let $A$ be a Hopf algebra over $k$.

\begin{definition}
    An $A$-module $M$ is called \emph{locally finite} if, for all $m \in M$,
    \[
    \dim_k(A \cdot m) < \infty.
    \]
    An $A^*$-module $M$ is called \emph{rational} if there exists $m_i \in M$ and $h_i \in A$, for all $m \in M$, such that
    \[
    f \cdot m = \sum f(h_i)m_i
    \] for all $f \in A^*$.
\end{definition}

For a locally finite $A$-module $M$, one can choose a basis $\{ v_i \}$ of $A \cdot m$ for all $m \in M$.
Then, we have $a \cdot v_j = c^i_j(a) v_i$ for some $c^i_j(a) \in A^\circ$, note that $c^i_j(1) = \delta^i_j$ and it is a matrix coefficient of $a$ w.r.t. the basis $\{v_i\}$.
Hence we might define a right coaction of $A^\circ$
\[
    \delta(m) := m^jv_i \otimes c^i_j
\]
where $m = m^j v_j$.
Conversely, for a $A^\circ$-comodule $M$, we define $A$-action on $M$ using the pairing between $A$ and $A^\circ$, and one can shows that any $A$-orbit of an elemnent of $M$ is of fintie dimensional. 
Hence a locally finite $A$-module is equivalent to a $A^\circ$ comodule.

For a rational $A^*$-module $M$, one may define a coaction of $A$ by $\delta(m) := \sum m_i \otimes h_i$, conversely, any $A$-comodule is a rational $A^*$-module.

Therefore, by the above arguments, we have the following propositon.

\begin{proposition}
    For a Hopf algebra $A$,
    \[
    \text{locally finite } A \text{-module} \Longleftrightarrow A^\circ\text{-comodule} \Longleftrightarrow \text{rational } (A^\circ)^*\text{-module}.
    \]
\end{proposition}
Note that each arrow reverse the direction of (co)action.

We may need the statement for $A^\circ$-module.

\begin{proposition}
    \(
    A^\circ\text{-module}
    \)
\end{proposition}


\subsection*{dual quasitriangularity of $D(H)$}.

Let $H$ be a Hopf algebra over a field $k$.
Then Drinfeld's quantum group $D(H)$ is uniquely determined by two sub Hopf algebra inclusions
\[
\imath: H \hookrightarrow D(H), \quad \jmath: H^{\circ, \text{op}} \hookrightarrow D(H)
\]
such that $\jmath(f)\imath(x) = \langle f_1, x_1 \rangle \langle f_3, S(x_3) \rangle \imath(x_2) \jmath(f_2)$ for all $f \in H^\circ, x \in H$.

% We first interpret $D(H)$ via YD-module over $H$:
% \begin{proposition}
%     \[
%     \mathcal{YD}^H_H \cong D(H) \mathrm{-Mod}_{H\text{-rat}}
%     \]
% \end{proposition}
% \begin{proof}
%     There is a 
% \end{proof}

% We verify the following proposition in \cite{Kashaev21}.

% \begin{proposition}
%     $D(H)^\circ$ has a dual quasitriangular structure with the $R$-form $\mathscr{R}: D(H)^\circ \otimes D(H)^\circ \rightarrow k$ defined by
%     \[
%     \mathscr{R}(x \otimes y) := \langle x, (\jmath \circ \imath^\circ)(y) \rangle
%     \]
% \end{proposition}
% \begin{proof}
    
% \end{proof}